\setcounter{section}{2}
\setcounter{theorem}{10}
\section{Characterization via Primes} \label{sec:characterization}

In this section we prove \autoref{thm:prime-characterization}, namely $\PPA_q = \amp_{p|q}\, \PPA_p$. The theorem follows by combining the following two ingredients.
\begin{itemize}[leftmargin=1.5cm,itemsep=3pt]
\item[\secref{sec:coprime}:]
$\PPA_{q r} = \PPA_q\bamp\PPA_r$ for any coprime $q$ and $r$.
\item[\secref{sec:prime-power}:]
$\PPA_{p^k} = \PPA_p$ for any prime power $p^k$.
\end{itemize}

\subsection{Coprime case} \label{sec:coprime}
\paragraph{\boldmath$\PPA_{qr} \supseteq \PPA_q \bamp\PPA_r$.}
We show that $\lonely_q \bamp\lonely_r$ reduces to $\lonely_{qr}$. Recall that an instance of $\lonely_q \bamp\lonely_r$ is a tuple $(C,V^*, b)$ where $(C,V^*)$ describes an instance of either $\lonely_q$ or $\lonely_r$ as chosen by $b\in\{0,1\}$. Suppose wlog that $b=0$, so the input encodes a $q$-dimensional matching $G = (V,E)$ over $V = [q]^n$ with designated vertices $V^*\subseteq V$, $|V^*| \not\equiv 0 \mod{q}$. We can construct a $qr$-dimensional matching $\overline{G} = (\overline{V},\overline{E})$ on vertices $\overline{V}\coloneqq V \times [r]$ as follows: For every hyperedge $e \coloneqq \set{v_1, \ldots, v_q} \in E$, we include the hyperedge $e \times [r]$ in $\overline{E}$. We let the designated vertices of $\overline{G}$ be $\overline{V}^* \coloneqq V^* \times [r]$. Note that $|\overline{V}^*| \not\equiv 0 \mod{qr}$. It is easy to see that a vertex $(v,i)$ is isolated in $G'$ iff $v$ is isolated in $G$. This completes the reduction since $\overline{V}$ is efficiently indexable, and the neighbors of any vertex in $\overline{V}$ are locally computable using black-box access to $C$. 

\paragraph{\boldmath $\PPA_{qr} \subseteq \PPA_q \bamp\PPA_r$.}
We show that $\bipartite_{qr}$ reduces to $\bipartite_q \bamp\bipartite_r$. Our input instance of $\bipartite_{qr}$ is a circuit $C : \bit^n \to (\bit^n)^k$ that encodes a bipartite graph $G = (V\cup U, E)$ with a designated node $v^* \in V$. If $\deg(v^*) \equiv 0 \mod{qr}$, then we already have solved the problem and no further reduction is necessary. Otherwise, if $\deg(v^*) \not\equiv 0 \mod{qr}$, we have, by the coprime-ness of $q$ and $r$, that either $\deg(v^*) \not\equiv 0 \mod{q}$ or $\deg(v^*) \not\equiv 0 \mod{r}$. In the first case (the second case is analogous), we can simply view $(G,v^*)$ as an instance of $\bipartite_q$, since vertices with degree $\not\equiv 0 \mod{q}$ in $G$ are also solutions to $\bipartite_{qr}$.

\subsection{Prime power case} \label{sec:prime-power}

$\PPA_{p^k} \supseteq \PPA_{p}$ follows immediately from our proof of $\PPA_{qr} \supseteq \PPA_q \bamp\PPA_r$, which didn't require that $q$ and $r$ be coprime. It remains to show $\PPA_{p^k} \subseteq \PPA_{p}$. We exploit the following easy fact.

\begin{fact}\label{fact:primes}
For all primes $p$, it holds that,\vspace*{-1mm}
\begin{align}
\text{for integers } t, c > 0: && \binom{c \cdot p^{t}}{p^t} \equiv 0 \mod{p} &\quad \text{if and only if }\ c \equiv 0 \mod{p}\label{eq:lucas-1}\\
\text{for integer } k > 0: && \binom{p^k}{i} \equiv 0 \mod{p} &\quad \text{for all }\ 0 < i < p^k\label{eq:lucas-2}
\end{align}

\end{fact}

\begin{figure}[t]
\centering
\begin{tikzpicture}
\tikzset{
	vert/.style = {circle, fill=black, minimum size=5pt, inner sep=0pt, outer sep=0pt},
	hedge/.style = {rectangle, rounded corners=4pt, above right, draw, minimum height=1.8cm, minimum width=1.8cm},
	medge/.style = {rectangle, rounded corners=4pt, draw, minimum height=0.4cm, minimum width=1.45cm, thick, right}
}
\node[vert] (A0) at (0,0) {};
\node[vert] (A1) at ([shift={(1,0)}]A0) {};
\node[vert] (A2) at ([shift={(1,1)}]A0) {};
\node[vert] (A3) at ([shift={(0,1)}]A0) {};
\node[hedge] at ([shift={(-0.4,-0.4)}]A0) {};
\node[medge, minimum width=1.45cm, draw=Gred] (A01) at ([shift={(-0.25,0)}]A0) {};
\node[medge, minimum width=1.45cm, draw=Gred] (A23)at ([shift={(-0.25,0)}]A3) {};
\draw[ultra thick, Gred] (A01) -- (A23);

\node[vert] (B0) at ([shift={(3.5,0)}]A0) {};
\node[vert] (B1) at ([shift={(1,0)}]B0) {};
\node[vert] (B2) at ([shift={(1,1)}]B0) {};
\node[vert] (B3) at ([shift={(0,1)}]B0) {};
\node[hedge] at ([shift={(-0.4,-0.4)}]B0) {};
\node[medge, minimum width=2.95cm, draw=Gblue] (B1C0) at ([shift={(-0.25,0)}]B1) {};
\node[medge, minimum width=2.95cm, draw=Gblue] (B2C3)at ([shift={(-0.25,0)}]B2) {};
\draw[ultra thick, Gblue] (B1C0) -- (B2C3);

\node[vert] (C0) at ([shift={(3.5,0)}]B0) {};
\node[vert] (C1) at ([shift={(1,0)}]C0) {};
\node[vert] (C2) at ([shift={(1,1)}]C0) {};
\node[vert] (C3) at ([shift={(0,1)}]C0) {};
\node[hedge] at ([shift={(-0.4,-0.4)}]C0) {};
\node[medge, minimum width=2.45cm, draw=Ggreen] (C1D0)at ([shift={(-0.25,0)}]C1) {};
\node[medge, minimum width=2.65cm, draw=Ggreen, rotate=-27] (C2D0)at ([shift={(-0.2,0.1)}]C2) {};
\draw[ultra thick, Ggreen, out=130, in=-100] ([shift={(0,0.12)}]C1.north) -- ([shift={(0.05,-0.15)}]C2.south);

\node[vert] (D0) at ([shift={(2,0)}]C1) {};
\node[vert] (D1) at ([shift={(2,0)}]D0) {};
\node[vert] (D2) at ([shift={(2,1)}]D0) {};
\node[vert] (D3) at ([shift={(0,1)}]D0) {};
\node[medge, minimum width=2.45cm, draw=Gyellow] (D1D2)at ([shift={(-0.25,0)}]D3) {};

\node[rectangle, rounded corners=4pt, above, draw, minimum height=1.8cm, minimum width=0.8cm, dashed] at ([shift={(0,-0.4)}]D1) {};
\node at ([shift={(0.9,0.5)}]D1) {$V^*$};

\end{tikzpicture}
\caption{Illustration of the proof of $\PPA_{p^k} \subseteq \PPA_{p}$ for $p=2$, $k=2$, $n=2$, $t=1$. In black, we indicate the $4$-dimensional matching $G$. In color, we highlight some of the vertices of $\overline{G}$ and the edges between them. The vertices of $\overline{G}$ in \textcolor{Gred}{red}, \textcolor{Gblue}{blue} and \textcolor{Ggreen}{green} are paired up and hence are non-solutions; whereas the vertex in \textcolor{Gyellow}{yellow} is isolated and not in $\overline{V}^*$ and hence a solution.}
\label{fig:ppa-prime-pow}
\end{figure}

\noindent We reduce $\lonely_{p^k}$ to $\lonely_p$. Our instance of $\lonely_{p^k}$ is $(C,V^*)$ where $C$ implicitly encodes a ${p^k}$-dimensional matching $G = (V = [p^k]^n, E)$ and a designated vertex set $V^* \subseteq V$ such that $|V^*| \not\equiv 0 \mod{p^k}$.

Let $p^t$, $0 \le t < k$, be the largest power of $p$ that divides $|V^*|$. Through local operations we construct a $p$-dimensional matching hypergraph $\overline{G} = (\overline{V}, \overline{E})$ over vertices $\overline{V} \coloneqq \binom{V}{p^t}$ (set of all size-$p^t$ subsets of $V$) with designated vertices $\overline{V}^*\coloneqq\binom{V^*}{p^t}$. From \eqnref{eq:lucas-1}, we get that $|\overline{V}| \equiv 0 \mod{p}$ and $|\overline{V}^*| \not\equiv 0 \mod{p}$.

We will describe an algorithm that on vertex $\overline{v} \in \overline{V}$ outputs a hyperedge of $p$ vertices that contains $\overline{v}$ (if any). To this end, first fix an algorithm that for any set $e \coloneqq \set{u_1, \ldots, u_{p^k}} \subseteq V$ and for any $1 \le i \le p^t$, computes some ``canonical'' partition of the set $\binom{e}{i}$ into subsets of size $p$, and moreover assigns a canonical cyclic order within each such subset. This is indeed possible because of \eqnref{eq:lucas-2}, since $t < k$.\\[-0.2cm]

\noindent Given a vertex $\overline{v} \coloneqq \set{v_1, \ldots, v_{p^t}} \in \overline{V}$,\vspace{-0.2cm}
\begin{itemize}[label=$\triangleright$,leftmargin=0.6cm]
\item Compute all edges $e_1, \ldots, e_\ell \in E$ that include some $v \in \overline{v}$.
\item For edge $e_j$, define $S_j \coloneqq e_j \cap \overline{v}$ and let $S_j^1, \ldots, S_j^{p-1}$ be the remaining subsets in the same partition as $S_j$ in the canonical partition of $\binom{e_j}{|S_j|}$, listed in the canonical cyclic order starting at $S_j$. Also, let $S_0$ be the set of untouched vertices in $\overline{v}$. Observe that $\overline{v} = S_0 \cup S_1 \cup \ldots \cup S_\ell$.
\item Output neighbors of $\overline{v}$ as the vertices $\overline{v}_1, \ldots, \overline{v}_{p-1}$ where $\overline{v}_i \coloneqq S_0 \cup S_1^i \cup \ldots \cup S_{\ell}^i$.
\end{itemize}

\noindent It is easy to see that $\overline{v}$ is isolated in $\overline{G}$ iff all $v \in \overline{v}$ are isolated in $G$. Moreover, any isolated vertex in $\overline{V} \smallsetminus \overline{V}^*$ contains at least one isolated vertex in $V \smallsetminus V^*$; and a non-isolated vertex in $\overline{V}^*$ contains at least one non-isolated vertex in $V^*$ (in fact $p^t$ many).

The edges of $\overline{G}$ can indeed be computed efficiently with just black-box access to $C$. In order to complete the reduction, we only need that $\overline{V}$ is efficiently indexable. This is indeed standard; see \cite[\S2.3]{r31} for a reference. See \autoref{fig:ppa-prime-pow} for an illustration of the proof.

\begin{remark}
Note that the size of the underlying graph blows up polynomially in our reduction. We do not know whether a reduction exists that avoids such a blow-up, although we suspect that the techniques of \cite{r6} can be used to show that some blow-up is necessary for black-box reductions.
\end{remark}


 
